%% Đề 9 -- de-giua-hoc-ky-2-toan-9-nam-2025-2026-truong-thcs-cao-loc-lang-son.pdf
\deheader{ĐỀ SỐ 9}%
  {Đề kiểm tra Giữa Học kỳ — Năm học 2025--2026}%
  {Phòng GD\&ĐT | Chỉ phần Đại số}
\addcontentsline{toc}{section}{Đề 9 — Đề kiểm tra Giữa HK — Năm học 2025--2026}

\begin{de_cau}{Câu 1}{2,0 điểm}
\textbf{Phần Trắc nghiệm Đại số (chọn đáp án đúng):}
\begin{enumerate}[label=\arabic*.]
\item Giá trị của hàm số $y = 2x^2$ tại $x = -3$ là:
\begin{multicols}{4}
A. $18$ \newline B. $-18$ \newline C. $12$ \newline D. $-12$
\end{multicols}
\item Phương trình nào sau đây là phương trình bậc hai một ẩn?
\begin{multicols}{4}
A. $2x + 1 = 0$ \newline B. $x^2 - 3x + 2 = 0$ \newline C. $x^3 - 1 = 0$ \newline D. $2x + y = 3$
\end{multicols}
\item Hệ số $a, b, c$ của phương trình $2x^2 - 5x + 3 = 0$ lần lượt là:
\begin{multicols}{4}
A. $2; -5; 3$ \newline B. $2; 5; 3$ \newline C. $-2; 5; 3$ \newline D. $2; -5; -3$
\end{multicols}
\item Biệt thức $\Delta$ của phương trình $x^2 - 4x + 3 = 0$ bằng:
\begin{multicols}{4}
A. $4$ \newline B. $-4$ \newline C. $16$ \newline D. $8$
\end{multicols}
\end{enumerate}
\end{de_cau}

\begin{loi_giai_cau}
\textbf{1.} Thay $x = -3$ vào $y = 2x^2 \implies y = 2(-3)^2 = 2 \times 9 = 18$. \textbf{Chọn A.}

\textbf{2.} Phương trình bậc hai một ẩn có dạng $ax^2 + bx + c = 0$ ($a \neq 0$), đó là phương trình $x^2 - 3x + 2 = 0$. \textbf{Chọn B.}

\textbf{3.} Phương trình $2x^2 - 5x + 3 = 0$ có $a = 2, b = -5, c = 3$. \textbf{Chọn A.}

\textbf{4.} $\Delta = b^2 - 4ac = (-4)^2 - 4 \cdot 1 \cdot 3 = 16 - 12 = 4$. \textbf{Chọn A.}
\end{loi_giai_cau}

\begin{de_cau}{Câu 2}{1,5 điểm}
Vẽ đồ thị của hàm số: $y = \dfrac{1}{2}x^2$.
\end{de_cau}

\begin{loi_giai_cau}
Tập xác định: $D = \mathbb{R}$.
Bảng giá trị:
\begin{center}
\begin{tabular}{|c|c|c|c|c|c|}
\hline
$x$ & $-2$ & $-1$ & $0$ & $1$ & $2$ \\ \hline
$y = \frac{1}{2}x^2$ & $2$ & $\frac{1}{2}$ & $0$ & $\frac{1}{2}$ & $2$ \\ \hline
\end{tabular}
\end{center}
Đồ thị là đường cong Parabol $(P)$ có đỉnh tại gốc tọa độ $O(0; 0)$, nhận trục $Oy$ làm trục đối xứng, bề lõm hướng lên trên và đi qua các điểm $A(-2; 2), B(-1; 1/2), O(0; 0), C(1; 1/2), D(2; 2)$.
\end{loi_giai_cau}

\begin{de_cau}{Câu 3}{1,5 điểm}
Giải các phương trình bậc hai sau:
\begin{enumerate}[label=\alph*)]
\item $2x^2 - 7x + 3 = 0$.
\item $x^2 - 6x + 9 = 0$.
\end{enumerate}
\end{de_cau}

\begin{loi_giai_cau}
\textbf{a)} $2x^2 - 7x + 3 = 0$:
$\Delta = (-7)^2 - 4 \cdot 2 \cdot 3 = 49 - 24 = 25 > 0$.
Phương trình có hai nghiệm phân biệt:
\[x_1 = \frac{7 + 5}{4} = 3;\qquad x_2 = \frac{7 - 5}{4} = \frac{1}{2}.\]
Vậy $\boxed{S = \left\{\dfrac{1}{2}; 3\right\}}$.

\textbf{b)} $x^2 - 6x + 9 = 0 \iff (x-3)^2 = 0 \iff x = 3$.
Vậy phương trình có nghiệm kép $\boxed{x = 3}$.
\end{loi_giai_cau}

\begin{de_cau}{Câu 4}{1,5 điểm}
Cho phương trình $x^2 - (m-2)x - m + 1 = 0$ ($m$ là tham số).
\begin{enumerate}[label=\alph*)]
\item Chứng minh phương trình luôn có nghiệm với mọi $m$.
\item Gọi $x_1, x_2$ là hai nghiệm của phương trình. Tìm $m$ để $x_1^2 + x_2^2 - 3x_1 x_2 = 1$.
\end{enumerate}
\end{de_cau}

\begin{loi_giai_cau}
\textbf{a)} Ta có:
\[\Delta = [-(m-2)]^2 - 4(1)(-m+1) = m^2 - 4m + 4 + 4m - 4 = m^2 \ge 0 \quad \forall m.\]
Do đó phương trình luôn có nghiệm với mọi giá trị của tham số $m$.

\textbf{b)} Theo định lý Vi-ét:
\[\begin{cases} x_1 + x_2 = m - 2 \\ x_1 x_2 = -m + 1 \end{cases}\]
Biến đổi hệ thức:
\[x_1^2 + x_2^2 - 3x_1 x_2 = (x_1+x_2)^2 - 5x_1 x_2 = (m-2)^2 - 5(-m+1) = m^2 - 4m + 4 + 5m - 5 = m^2 + m - 1.\]
Theo đề bài:
\[m^2 + m - 1 = 1 \iff m^2 + m - 2 = 0 \iff (m-1)(m+2) = 0.\]
Vậy $\boxed{m = 1}$ hoặc $\boxed{m = -2}$.
\end{loi_giai_cau}
